% !TEX root = ../journal_0421052099.tex
% ---------------------------------------------------------------------
% THE PAPER'S HEADLINE SECTION since the September 2026 reframe, and
% the one protected hardest under the page budget: it is cut less than
% any other section because it is the contribution.
%
% Delta is ablated-minus-reference, so POSITIVE MEANS REMOVING THE
% COMPONENT IMPROVED THE METRIC. Say that once, plainly, before the
% first number, or every figure below reads backwards.
%
% Report the planner asymmetry, never the pooled average -- pooling
% averages away the entire result. Only the WebQSP arm reaches
% significance; the CWQ arm is a DIRECTION, not an effect, and must not
% be called "marginally significant" (p = 0.088).
%
% The three nulls get the same length as the positive result. That is
% the methodological argument, not an apology.
%
% Power belongs HERE, not deferred to Discussion. check_paper_numbers
% recomputes every figure in sec:power and pins each to its own
% sentence; do not reword those without re-running the check.
%
% CUT 2026-09-23 (third pass) to ~980 words, from ~1,600. The
% implementation detail bounding the backtracking and scoring nulls
% moved to Appendix D (app:nulldetail); both nulls keep their deltas,
% their p-values, their token savings and their mechanism sentence
% here. NO NULL LOST A NUMBER THE BODY NEEDS.
% ---------------------------------------------------------------------

\section{Component-Level Attribution}\label{sec:attribution}

Four conditions each remove exactly one component from the unmodified
system, as the supplementary material specifies, and re-run the
identical questions: no planner, no backtracking, no claim verification,
and embedding-only scoring ($\alpha = 1$). Each is compared against the
reference on the same half-split, drawn as half of every hop stratum
rounded down, which gives $n = 200$ on WebQSP and $n = 198$ on
ComplexWebQuestions, with McNemar's test on the paired per-question
outcomes. The reference row is the unmodified system restricted to the
same questions, read from the main-run records rather than re-run, so it
is provably the system that produced \Cref{tab:main}. \textbf{Throughout
this section a delta is ablated minus reference, so a positive number
means removing the component \emph{improved} the metric.} The $p$-values
in \Cref{tab:ablation} are reported as computed; a correction for
multiple comparisons was not among the pre-specified decisions, and
applying one after seeing the results would be a choice made on the
results, so the four tests are read as descriptive.

\begin{table}[!tb]
  \begin{center}
        \caption{Component ablations. $\Delta$F1 is ablated minus reference:
      positive means the system got better without the component. The
      interval is a paired percentile bootstrap over questions
      ($10{,}000$ resamples); $p$ from McNemar's exact test on paired
      per-question correctness.}
    \label{tab:ablation}
    \scriptsize
    \setlength{\tabcolsep}{2.5pt}
    \begin{tabular}{|l|rlrr|rlrr|}
      \hline
      & \multicolumn{4}{c|}{\textbf{WebQSP} ($n=200$)}
      & \multicolumn{4}{c|}{\textbf{CWQ} ($n=198$)} \\
      \textbf{Removed} & $\Delta$F1 & $95\%$ CI & $p$ & Tokens
                       & $\Delta$F1 & $95\%$ CI & $p$ & Tokens \\
      \hline
      Planner        & $+0.083$ & $[+0.038, +0.130]$ & \textbf{0.006} & $-31\%$
                     & $-0.047$ & $[-0.108, +0.013]$ & 0.088 & $-21\%$ \\
      Backtracking   & $-0.013$ & $[-0.038, +0.010]$ & 0.727 & $-7\%$
                     & $0.000$  & $[-0.031, +0.030]$ & 1.000 & $-18\%$ \\
      Verification   & $+0.004$ & $[-0.006, +0.018]$ & 1.000 & $-2\%$
                     & $-0.009$ & $[-0.024, +0.000]$ & 1.000 & $-6\%$ \\
      Model scoring  & $-0.016$ & $[-0.060, +0.026]$ & 0.664 & $-25\%$
                     & $-0.008$ & $[-0.046, +0.030]$ & 0.481 & $-31\%$ \\
      \hline
    \end{tabular}
  \end{center}
\end{table}

\subsection{When Planning Hurts}\label{sec:planner}

One component produced a detectable effect, and it produced it in the
wrong direction. \textbf{Removing the planner improves WebQSP F1 by
$0.083$} ($p = 0.006$) while cutting token use by $31\%$. Twenty-one
questions were answered correctly only without the planner, against six
correct only with it, so the cheaper system is the more accurate one.

The mechanism is over-decomposition. WebQSP is predominantly a
single-hop benchmark: $256$ of its $400$ questions need one hop. Asked
to decompose a question with no compositional structure, the planner
manufactures sub-objectives that are not really separable, the explorer
then scores candidate edges against a fragment of the question rather
than the question, and it discards exactly the context that would have
identified the right relation. The planner prompt carries a single-hop
exemplar to show that not decomposing is a valid output, and that is
evidently not enough. The per-stratum breakdown locates the damage: on
the WebQSP half-split, removing the planner raises one-hop Hits@1 from
$0.75$ to $0.85$ and two-hop Hits@1 from $0.84$ to $0.89$.

On ComplexWebQuestions the direction reverses, as the mechanism
predicts. That benchmark is genuinely compositional, only $137$ of $400$
questions are single-hop, and removing the planner costs $0.047$ F1.
\textbf{That arm does not reach significance} ($p = 0.088$), and we
report it as a direction rather than an effect; we do not call it
marginally significant, because at this sample size that phrase would
assert more than the data carries. The strata show a coherent mechanism
rather than noise: without the planner, one-hop Hits@1 \emph{improves}
from $0.49$ to $0.54$, two-hop collapses from $0.48$ to $0.34$, and
three-hop-plus moves from $0.62$ to $0.54$. That is the pattern in which
decomposition helps genuine composition and hurts questions that do not
need it, and it is the pattern the hop-strata table in the supplementary
material shows from the other side. Pooling the two datasets would
average a significant improvement against a non-significant harm and
report approximately nothing. The asymmetry \emph{is} the result:
decomposition is not a general-purpose benefit but a mechanism whose
value depends on whether the question has compositional structure to
find.

\subsection{Three Components with No Detected Effect}\label{sec:nulls}

The remaining three conditions moved nothing detectable, and we report
them at the same length as the one that did. A field that reports only
what moved cannot attribute anything.

\textbf{Backtracking.} Removing it changes WebQSP F1 by $-0.013$ ($p =
0.727$) and ComplexWebQuestions F1 by $0.000$ ($p = 1.000$), while
saving $7\%$ and $18\%$ of tokens. The mechanism fires rarely: the
reference condition records $38$ and $108$ backtracks over the $398$
paired questions, and on those same runs the meter refuses a further
attempt on $28$ of them, $7.0\%$. A component this seldom exercised
cannot show a large effect. One implementation detail bounds the null
further and is recorded rather than smoothed over: the backtracker bans
the edges of the most recent pass but restores an often older snapshot,
so the null describes backtracking as implemented rather than as
specified (detailed in the supplementary material). The accurate summary
is that the search rarely needs to undo a decision, not that undoing is
worthless.

\textbf{Model scoring.} Setting $\alpha = 1$ removes the language model
from edge selection entirely. F1 falls by $0.016$ and $0.008$, neither
significant, while tokens drop $25\%$ and $31\%$. The model term is the
single most expensive part of the navigation loop, and this experiment
cannot establish that it earns its cost. The development-set sweep that
froze $\alpha$ points the same way at a smaller scale, and one confound
goes with the condition; The supplementary material gives both.

\textbf{Claim verification.} Removing the checking routes of the layer
AGR was designed around, while leaving the drafting call unchanged,
changes WebQSP F1 by $+0.004$ and ComplexWebQuestions F1 by $-0.009$, $p
= 1.000$ on both. Across the half-splits --- $200$ questions on WebQSP
and $198$ on ComplexWebQuestions --- the two conditions disagreed on
exactly one question each. \textbf{Claim verification does not
measurably change how often AGR is right.} The result is expected rather
than disappointing, for the reason \Cref{sec:verification} gave before
any of these numbers appeared: on most questions the layer is a
pass-through certifying a draft that was already grounded. On test it
asks for a repair on $52$ of $800$ questions, exploration actually
resumes on $28$, and $13$ end grounded that did not start so. A
mechanism that alters the output on a few per cent of questions cannot
move an aggregate accuracy metric, and an ablation is the right
instrument for establishing that rather than assuming it either way.
What the layer delivers is therefore not an accuracy gain but the output
contract: an answer paired with the traversed triples that support it,
and an entity list that is a deterministic function of the verdicts
rather than a second generation. That is an auditability claim, not an
accuracy claim, and \Cref{sec:limits-contribution} states how far it
reaches; \Cref{sec:groundedness} has already shown that the layer cannot
claim credit for groundedness.

\subsection{Statistical Power}\label{sec:power}

These are nulls at low power, not demonstrations of absence. McNemar's
exact test sees only the \emph{discordant} pairs, and the null
conditions produced very few. Removing claim verification changed the
outcome on exactly one question per dataset, and with a single
discordant pair no split whatsoever can reach significance, so that row
is not a test that failed to reject: it is the two conditions agreeing
on $396$ of the $398$ paired questions, a statement about how rarely the
layer alters an answer rather than about statistical power. Backtracking
produced $8$ and $11$ discordant pairs, and model scoring $21$ and $18$.
At those counts the smallest imbalance the exact test could have called
significant is $8$ and $9$ questions for backtracking, meaning all eight
would have had to fall the same way, and $11$ and $10$ for model
scoring, which is between $4.0$ and $5.5$ points of accuracy.
Differences below that were undetectable at any power, whatever their
true size.

For calibration: detecting a true $2{:}1$ discordance ratio at $80\%$
power requires roughly $72$ discordant pairs, and the largest of these
conditions produced $21$, at which count the smallest ratio detectable
at $80\%$ power is nearer $4{:}1$, so a true $2{:}1$ effect would have
been caught about a quarter of the time. \textbf{Three of four
components show no effect at this power}; none is shown to have no
effect. One more count bounds what any single removal could have done:
the union of the four discordant sets holds $42$ of the $200$ WebQSP
questions and $55$ of the $198$ on ComplexWebQuestions. On the remaining
$158$ and $143$, no condition changed the outcome. The intervals in
\Cref{tab:ablation} say the same thing in F1, since the planner's WebQSP
interval is the only one of the eight that does not contain zero. Each
condition was also run once and the components were removed one at a
time, so three nulls are consistent with the three being individually
dispensable and jointly necessary, and equally with a system that has
none of them matching this one at a fraction of the tokens; the
condition that separates those readings was not run, and
\Cref{sec:outlook} names it.
